\chapter{Оценка производительность узлов сети Интернет в режиме реального времени}\label{ch:ch4}

\section{Предобработка данных о производительности сети}\label{sec:ch4/sect1}

Для сравнительного анализа и апробации методов оценки производительности сетевой структуры была проведена серия измерений доступности и времени отклика портала государственных услуг Российской Федерации — gosuslugi.ru. Портал был выбран в качестве объекта исследований в виду его государственной важности, что снижает риски проблем с оборудованием на стороне конечного адреса. Эксперимент выполнен с использованием распределённой измерительной платформы, описанной в предыдущих разделах.

Наблюдения проводились в течение 16 суток — с 26 мая 2024 года, 18:32:00 по UTC, по 11 июня 2024 года, 18:32:00 по UTC. Интервал между измерениями составлял 1200 секунд (или же 20 минут), что обеспечило равномерное временное покрытие в течение всего периода мониторинга.

В ходе эксперимента было получено свыше 550 000 индивидуальных измерений времени ответа, собранных со всех узлов. Географическое распределение измерительных устройств неравномерно: значительная часть зондов (более 240) сосредоточена в центральном регионе, что отражено на карте размещения измерительных точек (Рисунок \ref{fig:map}). Данная концентрация может оказывать влияние на пространственную репрезентативность результатов, однако обеспечивает высокую детализацию в наиболее населённом и социально активном регионе страны.

В измерительном процессе участвовали все доступные зонды RIPE Atlas, размещённые на территории РФ — всего 561 устройство. Каждая станция выполняла периодическое измерение времени отклика путём отправки четырёх ICMP-запросов c объёмом пакета в 48 байт. В качестве результата измерения фиксировалось среднее значение времени отклика по четырём запросам, что позволило снизить влияние кратковременных флуктуаций и повысить устойчивость оценки.

Полученные данные легли в основу последующего анализа доступности, устойчивости и временной динамики производительности сервиса в условиях реальной эксплуатации.


\begin{figure}[h]
	\centerline{\includegraphics[width=0.8\linewidth]{time_series/map}}
	\caption{Места расположения зондов, задействованных в проведенных измерениях }
	\label{fig:map}
\end{figure}


Для обеспечения корректности последующего анализа потребовалась предварительная обработка данных, обусловленная наличием в исходной выборке ряда  особенностей. Одной из таких особенностей является нестационарность временного шага между измерениями, вызванная стохастическими процессами, неизбежными в условиях реального функционирования измерительной инфраструктуры. В частности, для зонда с идентификатором prb\_id = 11337 было зафиксировано изменение интервала между последовательными измерениями в диапазоне от 1187 до 1215 секунд (см. Рисунок \ref{fig:hist_time}). Указанный разброс требует применения соответствующих методов синхронизации и интерполяции временных рядов для минимизации погрешностей при последующем анализе.

В ходе анализа полученных данных было отмечено несколько особенностей в данных: ранее упомянутый разброс в измерениях, а также неполнота (нервное кол-во полученных измерений отклика у различных устройств). Поэтому для получения точного и репрезентативного анализа была проведена следующая предобработка данных:

\begin{itemize}
	\item Нумерация полученных измерений на всех устройствах с учетом фактического разброса шага.
	\item Фильтрация устройств: для дальнейшего анализа были оставлены лишь те устройства, которые проработали 95\% от всего периода измерений. В результате чего, в дальнейшем анализе использовались 505 устройств (Неполнота некоторых точек измерений отображена на Рисунке \ref{fig:prb_count})
	\item Заполнение пропусков: в случае пропущенного измерения значения устанавливались медианным значением для всего периода.
\end{itemize}

Итоговый состав использованных в дальнейшем устройств по типу и версии платформы отображен в Таблице \ref{tab:ping_table}.


\begin{figure}[htbp]
	\centering
	\includegraphics[width=0.5\linewidth]{time_series/hist_time}
	\caption{Разброс времени между измерениями на примере устройства id\_probe = 12794}
	\label{fig:hist_time}
\end{figure}


\begin{figure}[htbp]
	\centering
	\includegraphics[width=0.7\linewidth]{time_series/prb_counts}
	\caption{Распределение количества успешно проведенных замеров на устройстве}
	\label{fig:prb_count}
\end{figure}


\begin{table}[!ht]
	\caption{Распределение устройств по типу и версии в предобработанной выборке}
	\centering
	\begin{tabular}{|c|c|}
		\hline
		Тип/Версия & Кол-во устройств \\ \hline
		system: V3 & 187  \\ \hline
		system: V4 & 125  \\ \hline
		Software   & 116  \\ \hline
		system: V2 &  43  \\ \hline
		Anchor     &  20  \\ \hline
		system: V1 &  14  \\ \hline
	\end{tabular}
	\label{tab:ping_table}
\end{table}

\section{Применение методов кластеризации на данных о производительности сети}\label{sec:ch4/sect2}

%В данном разделе представлены результаты измерения времени отклика в различных регионах страны с целью выявить глобальные и региональные особенности работы сети. Для измерения времени отклика использовалась система RIPE Atlas, особенности работы которой описаны в работах \cite{ripe} и \cite{RIPEanalisys}.  В данной системе есть подходящий для задачи проводимого исследования режим измерений, обозначаемый как «Ping». В данном режиме выбранные станции отправляют запрос на введенный пользователем адрес в соответствии с выставленными параметрами, после чего фиксируется время отклика. Результаты измерений записываются в хранилище системы RIPE Atlas и доступны к выгрузке по окончанию измерений. Несмотря на критическую важность Интернет-инфраструктуры, единых и массовых систем мониторинга не существует. В свою очередь, RIPE Atlas, будучи некоммерческим проектом от энтузиастов, объединяет исследователей по всему миру и позволяет проводить измерения благодаря массовому внедрению и открытому доступу.

В целях детального изучения свойств полученных данных был выполнен кластерный анализ временных рядов. Для этого использовался алгоритм K-средних (K-means), а для вычисления расстояния между временными рядами использовался алгоритм динамической трансформации временной шкалы (Dynamic Time Warping - DTW) \cite{dtw}. Также стоит отметить, что результаты такой кластеризации могут помочь в улучшении эффективности моделей предсказания временных рядов в дальнейшей работе.

Для подбора оптимального кол-ва параметров кластеризации результаты сравнивались между собой по среднему коэффициенту силуэта, который базируется на расстоянии между кластерами и внутри кластеров \eqref{sil}.

\begin{equation}
	\label{sil}
	s = \frac{b - a}{max(a, b)} ,
\end{equation}
где a – среднее расстояние от рассматриваемого объекта до остальных объектов того же кластера, b – среднее расстояние от рассматриваемого объекта до объектов ближайшего кластера.  Исходя из данного критерия оптимальным значением для количества кластеров является $n=3$ (Рисунок \ref{fig:silh}).



\begin{figure}[htbp]
	\centering
	\includegraphics[scale=0.55]{time_series/silh}
	\caption{Значения коэфициента силуэта при применении алгоритма кластеризации для различного выставленного количества кластеров}
	\label{fig:silh}
\end{figure}


В ходе применения кластеризации с найденными оптимальными параметрами были получены 3 кластера, в которых можно выделить три различных паттерна поведения в производительности узлов сети (Рисунок \ref{fig:clusters}). Соотношение кластеров в полученных данных отображено в таблице \ref{tab:cluster_balance}. Из-за наличия различных семейств временных рядов можно сделать вывод о необходимости сегментации временных рядов при смене паттерна, а также обучении разных моделей на разных кластерах и сегментах временных рядов для увеличения эффективности в задаче предсказания временных рядов.


\begin{figure}[htbp]
	\centering
	\includegraphics[width=0.9\linewidth]{time_series/clusters}
	\caption{Визуальное отображение применения кластеризации к временным рядам времени отклика. Черным отображены полученные ядра кластеров, а серым - множество временных рядов, отнесенных к этому кластеру}
	\label{fig:clusters}
\end{figure}



\begin{table}[h!]
	\caption{Соотношение детектируемых кластеров в используемых зондах.}
	\centering
	\begin{tabular}{|c|c|}
		\hline
		Кластер     & Доля в данных \\ \hline
		Кластер №1  & 68 \%  	    \\ \hline
		Кластер №2  & 13 \% 	    \\ \hline
		Кластер №3  & 19 \%  	    \\ \hline
	\end{tabular}
	\label{tab:cluster_balance}
\end{table}



\section{Методы оценки и прогнозирования быстродействия сети }\label{subsec:ch4/sect3}

В настоящем разделе рассматриваются методологические аспекты анализа и оценки производительности сетевой инфраструктуры на территории страны, с акцентом на количественную оценку показателей быстродействия Интернета. Представлен комплекс методов анализа, включающий обработку аномальных значений, сравнительный анализ географически распределённых зон, а также декомпозицию временных рядов с целью выявления структурных компонент (таких как тренд, сезонность и шум). 

Особое внимание уделено факторам, оказывающим значимое влияние на качество сетевого соединения, к которым относятся: географическое расстояние до целевых серверов, временные характеристики измерений (включая суточную динамику и фазы активности), а также принадлежность к конкретному интернет-провайдеру. Анализ этих факторов позволяет выявить ключевые детерминанты сетевой производительности и оценить их вклад в общую вариативность измеряемых показателей.

Здесь представлены результаты анализа времени отклика, которое было получено с помощью системы RIPE Atlas. Подробнее о времени и характере проводимых измерений можно повторно ознакомиться в разделе \ref{sec:ch4/sect1}


\subsection{Анализ зависимости времени отклика от расстояния}

Для анализа зависимости времени отклика от расположения станции относительно сервера на данном этапе исследования были использованы две модели: линейная регрессия с регуляризацией и метод опорных векторов c линейным ядром (SVR). В качестве признаков использовались полученные из RIPE Atlas широта и долгота места установки точки измерения, а также расстояние до исследуемых серверов исходя из координат конечного адреса, полученных через сервис geolocation.com.

Согласно результатам моделирования лучше всего показала себя линейная регрессия с регуляризацией типа Lasso (Рисунок \ref{fig:linreg1}, \ref{fig:linreg2}), показывая на кроссвалидации медианное значение $R^2$ более 0,84, при этом веса модели для широты и долготы оказываются зануленными. Подобранные оптимальные параметры указаны в Таблице \ref{tab:params}. Стоит отметить что данное моделирование проводилось только для аппаратных устройств V3 в силу их минимальных искажений результатов \cite{ripe, RIPEanalisys}. Для использования данных со всех устройств необходимо проводить калибровку полученных показателей отклика.

\begin{figure}[ht]
	\centerfloat{
		\hfill
		\subcaptionbox[List-of-Figures entry]{Подбор модели \label{fig:linreg1}}{%
			\includegraphics[width=0.49\linewidth]{time_series/crossval}}
		\hfill
		\subcaptionbox{Применение оптимального варианта \label{fig:linreg2}}{%
			\includegraphics[width=0.49\linewidth]{time_series/lasso}}
		\hfill
	}
	\caption{Результаты анализа зависимости времени отклика от расстояния. На рисунке а) отображены диаграммы размаха для коэффициента $R^2$ для обученных регрессионных моделей зависимости медианного времени отклика от географического расположения станций. SVR – регрессия методом опорных векторов, Lasso и Ridge – линейные регрессии с регуляризацией соответствующего типа, LinReg – линейная регрессия без регуляризации. На рисунке б) - визуализация результатов применения линейной регрессии с регуляризацией типа Lasso}\label{fig:linreg}
\end{figure}




\begin{table}[h!]
	\caption{Оптимальные гиперпараметры для моделей, использованных для оценки зависимости rtt от расстояния между узлами}
	\centering
	\begin{tabular}{|c|c|c|}
		\hline
		Модель & Регуляризация & Подобранные гиперпараметры  \\ \hline
		Линейная регрессия & Lasso & $\alpha$ = 27,825  \\ \hline
		Линейная регрессия  & Ridge & $\alpha$ = 12915,49  \\ \hline
		SVR & - & C = 0,01; $\epsilon$= 1,0  \\ \hline
	\end{tabular}
	\label{tab:params}
\end{table}


Среди анализируемых устройств выявлено более 9\% зондов, для которых наблюдаются значимые отклонения наблюдаемых значений времени отклика от предсказанных значений, полученных с использованием наиболее оптимальной модели среди ранее рассмотренных (отображены на Рисунке \ref{fig:anomalies}). Эти отклонения превышают 1,5 интерквартильного размаха (IQR) относительно первого и третьего квартилей распределения остатков модели( разность между наблюдаемыми и предсказанными значениями rtt). Анализ характера отклонений показывает, что они носят односторонний характер и преимущественно соответствуют увеличению времени отклика по сравнению с модельными оценками. Такая систематическая положительное смещение указывает на потенциальные проблемы с производительностью оборудования или сетевой инфраструктуры в соответствующих локациях. Полученные результаты свидетельствуют о необходимости проведения детального аудита и оптимизации функционирования сетевых компонентов в выявленных участках. Дополнительные сведения о характеристиках аномальных зондов представлены в Таблице \ref{tab:anomalies} Приложения \ref{app:B}.


\begin{figure}[ht]
	\centerfloat{
		\hfill
		\subcaptionbox[List-of-Figures entry]{\label{fig:anomalies}}{%
			\includegraphics[width=0.49\linewidth]{time_series/anomaly}}
		\hfill
		\subcaptionbox{\label{fig:outputs}}{%
			\includegraphics[width=0.49\linewidth]{time_series/tsv2}}
		\hfill
	}
	\caption{ а) - наблюдаемые данные и оценка модели (прямая). Треугольными метками изображены <<аномальные>> измерения. б) - исходный вид полученных данных до предобработки на примере устройства prb\_id=11337. Временной шаг между отчетами 20 мин }\label{fig:anomalies_outputs}
\end{figure}




\subsection{Анализ полученных временных рядов}

Следующим этапом проведенного исследования полученных данных является анализ данных каждого устройства с точки зрения временных рядов: декомпозиция (выявление трендовых, сезонных и цикличных компонент), изучение автокорреляции исходных и продифференцированных рядов, проверка стационарности. Перед этим данные были очищены от выбросов (пример таковых можно увидеть на Рисунок \ref{fig:outputs}). Для этого из каждого набора данных для всех станций время отклика, выходящее за 1,5 интерквартильных размаха от 1 и 3 квартиля были заменены на новые значения c использованием линейной интерполяции по соседним точкам.

Далее для анализа общей картины полученных временных рядов был использован следующий подход: для каждого временного ряда вычислялось значение частичной автокорреляции (PACF) для значений шага от 1 до 100, далее вычислялась медиана PACF для каждого значения шага.  PACF определяется следующим образом:

\begin{equation}
	PACF(k)=\left\{\begin{array}{ccccccccccc}
		corr(y_{t+k}, y_t) , k=1,\\
		corr(y_{t+k} - y_{t+k}^{k-1}, y_t - y_t^{k-1}),k>1,
	\end{array}
	\right.
\end{equation}
где $y_t^{k-1}$ - линейная регрессия на $y_{t+1}, y_{t+2}, \dots , y_{t+k-1}$, минимизирующая среднеквадратичное отклонение (MSE). 

В полученных выше описанным путем данных не наблюдается характерной картины нестационарности исследуемых временных рядов (Рисунок \ref{fig:median}), а также тест стационарности Филлпса-Перрона указывает на стационарность (p-value << 0,05).

\begin{figure}[h] \centering
	\includegraphics[height=0.2\textheight]{time_series/pacf1}
	\caption{График медианной PACF после предобработки данных}
	\label{fig:median}
\end{figure}


Однако стоит отметить, что в данных присутствует высокочастотный шум. Для снижения его влияния был использован фильтр Савицкого-Галлея. Фильтр применялся с параметрами: окно w=20, степень полинома p=2. Результат применения такого фильтра можно увидеть на примере устройства с prb\_id = 12794 (Рисунок \ref{fig:filtered}). Сравнение разных параметров фильтра можно найти в Приложении.


\begin{figure}[h] \centering
	\includegraphics[height=0.2\textheight]{time_series/filter}
	\caption{График измерения времени отклика с течением времени для устройства prb\_id = 12794 до и после применения фильтра Савицого-Галлея. Временной шаг между отсчётами $\Delta t$ = 20 мин}
	\label{fig:filtered}
\end{figure}


Можно отметить, что присутствует сильная взаимосвязь с последним значением: для шага l=1 наблюдаем существенный всплеск значения автокорреляции (Рисунок \ref{fig:pacf1}).  После однократного дифференцирования с шагом 1 данный аспект устраняется и остаются только компоненты с шагом 20 (Рисунок \ref{fig:pacf2}).


\begin{figure}[ht]
	\begin{minipage}[t]{0.49\linewidth}
		\center{\includegraphics[height=0.2\textheight]{time_series/pacf2v2}}
		\caption{График медианы частичной автокорреляции (PACF) для полученных данных после применения фильтра Савицкого-Галея}
		\label{fig:pacf1}
	\end{minipage}
	\hfill
	\begin{minipage}[t]{0.49\linewidth}
		\center{\includegraphics[height=0.2\textheight]{time_series/pacf3v2}}
		\caption{График медианы частичной автокорреляции (PACF) для полученных данных после применения фильтра Савицкого-Галея и однократного дифференцирования ряда с шагом d=1}
		\label{fig:pacf2}
	\end{minipage}
\end{figure}


Таким образом, можно отметить, что после фильтрации и однократного дифференцирования в полученных временных рядах наблюдается сезонная компонента с периодом T = 20 (400 мин. или 6 часов 40 минут). Данные результаты дают понимание того, какие значения параметров следует выставить при использовании моделей, базирующихся на STL разложении, например, таких как автокорреляционные модели ARIMA \cite{arima}. 



\subsection{Спектральный анализ полученных значений RTT}

Для выявления и анализа сезонных компонент в полученных временных рядах был выполнен спектральный анализ на основе дискретного преобразования Фурье (ДПФ), а именно с использованием алгоритма быстрого преобразования Фурье (FFT) \cite{fft}. Полученные спектрограммы позволили оценить распределение мощности сигнала по частотному диапазону и выявить главные периодические составляющие. С целью обобщения спектральных характеристик по всем измерительным устройствам, индивидуальные спектрограммы были агрегированы. Результаты представлены в виде диаграмм размаха (или же <<ящиков с усами>>), отражающих распределение спектральной мощности на различных частотах (Рисунок \ref{fig:spectre1}), а также усреднённой спектральной мощности, демонстрирующей доминирующую частотную структуру во времени отклика (Рисунок \ref{fig:spectre2}).


В ходе анализа полученных результатов было выделено наличие следующей особенности в наблюдаемом спектре: компонента с периодом равном одним суткам (SNR=25,6 dB для медианного спектра). Также нельзя исключить наличие компонент с большим периодом (несколько дней, неделя и более), но имеющиеся данные не позволяют определить это с достаточной точностью. Для этих целей в ходе дальнейших исследований запланировано проведение измерений большей длительности. 



\begin{figure}[ht]
	\begin{minipage}[t]{0.49\linewidth}
		\center{\includegraphics[height=0.2\textheight]{time_series/spectre1v2}}
		\caption{Диаграммы размаха для спектрограмм времени отклика с рассматриваемых устройств}
		\label{fig:spectre1}
	\end{minipage}
	\hfill
	\begin{minipage}[t]{0.49\linewidth}
		\center{\includegraphics[height=0.2\textheight]{time_series/spectre2v2}}
		\caption{Медианная спектральная картина для времени отклика для рассматриваемых устройств для значений периода, выраженного в сутках}
		\label{fig:spectre2}
	\end{minipage}
\end{figure}


Также был проведен сравнительный фазовый анализ сезонных компонент во времени отклика для выявления наличия или отсутствия взаимосвязи синхронности в периодических изменениях быстродействия сети между устройствами в разных регионах. Для этого на имеющихся данных было применено преобразование Гильберта. Согласно результатам применения данного преобразования и последующего попарного сравнения временных рядов с разных устройств, средняя разница фаз составляет менее 0,05 радиан, что указывает на синхронность периодических изменений, а также независимость этих компонент от часового пояса устройства. Этот вывод подтверждается также при визуальном сравнении данных с 3 устройств: из Москвы, Омска и Владивостока (Рисунок \ref{fig:noremed_ts}). 


\begin{figure}[h] \centering
	\includegraphics[height=0.2\textheight]{time_series/phase1}
	\caption{Отнормированные показатели времени отклика для трех случайных устройств из трех городов. Время между отсчётами $\Delta t$ = 20 мин}
	\label{fig:noremed_ts}
\end{figure}

Данная особенность становится более наглядной после использования полосового фильтра Баттерворта \cite{ibm} с частотными параметрами lowcut = 1/78, highcut = 1/68 (в 1/<временной отсчет>, то есть в 1/20 min), которые соответствуют периодам в ~22.7 и 26 часов (Рисунок \ref{fig:filtered_ts}).


\begin{figure}[h] \centering
	\includegraphics[height=0.2\textheight]{time_series/phase2}
	\caption{Результат применения полосового фильтра Баттерворта на данных с трех случайных устройств из трех регионов}
	\label{fig:filtered_ts}
\end{figure}


Медианное время (по UTC) пика нагрузки исходя из анализа суточной сезонной компоненты наблюдаемых данных представлено в Таблице \ref{tab:cities}.

\begin{table}[h!]
	\caption{Медианное время пиковой нагрузки для трех устройств в разных городах}
	\centering
	\begin{tabular}{|c|c|}
		\hline
		Город & Время пиковой нагрузки  \\ \hline
		Москва & 12:17:03  \\ \hline
		Омск & 11:17:26  \\ \hline
		Владивосток & 11:35:25  \\ \hline
	\end{tabular}
	\label{tab:cities}
\end{table}


Время для пиковой нагрузки для трех рассматриваемых станций отличается, однако отличается оно меньше, чем на разницу часовых поясов. Таким образом можно сделать вывод, что время суточной пиковой нагрузки (именно пик суточной периодической составляющей) значимо не зависит от часового пояса, в котором расположено устройство. Таким образом эти колебания носят глобальный характер для всей сети страны.

%- спектральный анализ \\
%- подход с использование автокорреляционных функций\\
%- анализ кластеров ?\\
%- итоги\\
\section{Анализ вычислительной сложности}\label{subsec:ch4/sect4}

Проведем анализ вычислительной сложности разработанного подхода и аналогов для схожих задач. В анализе будем использовать асимптотические оценки (в виде <<О-большое>>), а также FLOPs. FLOPs (Floating Point Operations) — это число операций с плавающей точкой, необходимых для выполнения одного прохода модели машинного обучения. Эта метрика отражает вычислительную сложность и используется для оценки эффективности и ресурсоёмкости алгоритма. Также для краткости  операция одного прохода модели для единичного вычисления (т.е. на один временной отсчет вперед) в дальнейшем будет обозначаться как <<инференс>>.

%Перед оценкой сложности разработанных подходов к анализу времени отклика введем следующие обозначения:

%\begin{itemize}
%	%\item $K$ - количество кластеров
%	\item $T$- длинна временного ряда
%	\item $N$ - кол-во временных рядов, размер обучающий выборки
%	\item $M$ - число итераций обучения (например кол-во эпох)
%\end{itemize}


%Для начала приведем анализ сложности обучения для \textbf{K-means}. В ходе обучения необходимо рассчитать расстояние для каждого элемента до всех цетроидов кластеров (в начале они выбираются случайным образом). С учетом вычислительной сложности метрики DTW - $O(T^2)$ получаем итоговую сложность:

%\begin{equation}
%	O(N K T^2)
%\end{equation}

%Далее следует поиск минимума расстояния для каждого элемента из N:

%\begin{equation}
%	O(N K)
%\end{equation}

%После чего идет пересчет центроидов кластеров:

%\begin{equation}
%	O(N T^2)
%\end{equation}

%В итоге получаем, что сложность обучения будет:

%\begin{equation}
%	O(N K T^2 + NK + NT^2) = O(N K T^2)
%\end{equation}

%Также если учесть количество итераций алгоритма - M:

%\begin{equation}
%	O(M N K T^2)
%\end{equation}

%Сложность оценки кластера для нового элемента (инференса) состоит из вычисления метрики до каждого из центроидов, то есть:

%\begin{equation}
%	O(K T^2)
%\end{equation}

Рассмотрим сложность алгоритма \textbf{ARIMA}. Обозначим суммарное кол-во компонент AR и MA (а также их сезонные версии) как $F = p + q + P + Q$. 
%Тогда с учетом количества итераций M сложность обучения составит:

%\begin{equation}
%	O(M T F^2)
%\end{equation}

Тогда сложность же однократного инференса - $O(F)$. Если же переводить в количество операций с плавающей точкой, то учитывая наличие сложения и умножения для каждой компоненты получим - $2 \times F$ FLOPs.


По итогам анализа времени отклика была подобрана модель для неаномальных временных рядов. Для всех рядов наиболее оптимальными параметрами оказались: $p=1$, $d=0$, $q=1$, $P=1$, $Q=1$, а также в виду наличия сезонности - $D=1$. В ходе оценки данной модели использовалась отложенная выборка в 20\% от полученных рядов (пример этого отображен на Рисунке \ref{fig:arima162}). В итоге обучения такой модели и валидации на отложенной выборке были получены средние  значения метрик ошибок:

\begin{itemize}
	\item MAE = 0,1178
	\item RMSE = 0,1493
	\item MSE = 0,0362
	\item MAPE = 2,4573
\end{itemize}

Для проведения дальнейшей оценки сложности модели введем обозначение для горизонта предсказания - $T$. Это количество шагов, на которое модель дает оценку времени отклика. Оценки всех моделей будут проводится для  $T=100$. Тогда сложность для рассматриваемого подхода ARIMA составит  есть сложность можно оценить как $O(F \times T)$, тогда оценка количества операций составляет: $2 \times F \times T$  = 800 FLOPs.


\begin{figure}[h] \centering
	\includegraphics[height=0.3\textheight]{ARIMA/fitted_arima_162}
	\caption{Пример применения исследуемого подхода на примере устройства с идентификатором prb\_id=162. Train - обучающая выборка, Test - отложенная, Forecast - результат оценки модели}
	\label{fig:arima162}
\end{figure}

%Для дальнейшего сравнительного анализа рассмотрим одну из архитектур нейросетей для работы с временными рядами -LSTM. Добавим новое обозначение: H - кол-во слоев, B - размер батча. Тогда обучение за все эпохи выйдет:

%\begin{equation}
%	O(M N T H^2) 
%\end{equation}

Далее рассмотрим модели из современных работ, решающих схожие задачи анализа сети \cite{ge2022gnn, lstm_trans}. Как ранее упоминалось - часто в рассматриваемых задачах используются нейросетевые подходы, в частности в данных работах рассматриваются 3 архитектуры: долгая краткосрочная память (Long short-term memory; LSTM), Трансформеры (Transformers) и графовая нейронная сеть (Graph Neural Network, GNN).

Итак в работе \cite{lstm_trans} была использована нейросетевая модель \textbf{LSTM} со следующими параметрами:

\begin{itemize}
	\item Архитектура:  многошаговая, многослойная LSTM
	\item Количество слоев: 2 последовательно соединённых слоя
	\item Размерность скрытого состояния: $S=16$
	
\end{itemize}

Также стоит отметить, что модель включает в себя и дополнительные параметры, такие как остаточные соединения (residual connections), регуляризация методом <<дропаут>> (dropout), а также слои токенизации (разбиение временного ряда на последовательности фиксированной длины, т.н. "токены") для обеспечения устойчивого обучения.


На вход подается последовательность из $H$ токенов размерности $S$. Тогда сложность для одного шага: $O(S \times H \times N)$, где N — число слоёв (в данной работе это 2). Тогда сложность на всю последовательность: $O(H \times S \times S \times N)$. Модель делает оценку на $T$ шагов вперед. Сложность такой операции составляет:

\begin{equation}
	O(H \times S^2 \times N) + O(T \times S^2 \times N) = O((H+T) \times S^2 \times N)
\end{equation}

Оценим количество операций. С учетом наличие 4-х гейтов, входного и скрытого состояния на один инференс получается: $ 2 \times 2 \times [2 \times (1+1) \times 16^2  \times 2] = 8192$ FLOPs. Тогда оценка количества операций для 100 отсчетов составит 819\,200 FLOPs.

Также авторы \cite{lstm_trans} использовали \textbf{Transformers} со следующими параметрами:

\begin{itemize}
	\item Архитектура: стандартный энкодер-декодер с многообразным механизмом самовнимания (multi-head self-attention)
	\item Количество слоев: 6 слоёв энкодера и 6 слоёв декодера
	\item Количество блоков внимания: 4 головы на слой
	\item Увеличение размерности в полносвязных слоях: в 32 раза
	\item Скрытое пространство (размерность S): 16
	
\end{itemize}

Вторая модель также включается в себя <<дропаут>> на уровне 0,2, в качестве функции активации авторы использовали GELU. В итоге сеть имела около 78 тысяч весов.


Итак, оценим вычислительную сложность данной модели. На вход подается $H$ токенов на энкодер, и $L$ токенов декодер. Тогда с учетом блока self-attention: $O(H^2 \times S)$ на энкодер, $O(T^2 \times S)$ на декодер. Механизм cross-attention дает сложность $O(T \times H \times S)$. Далее слои feed-forward: $O(H \times S \times d_{ff})$ на энкодер, $O(T \times S \times d_{ff})$ на декодер (где $d_{ff} = 32 \times S = 512$). В итоге, с учетом того, что в данной работе $N_{enc} = N_{dec}$ получаем:

\begin{equation}
	\begin{split}
		O(N_{enc} \times (H^2 \times S + H \times S \times d_{ff})) + \\ O(N_{dec} \times (T^2 \times S + T \times H \times S + L \times S \times d_{ff})) = \\
		O(N(H^2+T^2)S + N(H+L)S d_{ff}) 
	\end{split}
\end{equation}


Оценим количество операций: на self-attention энкодер: $6 \times (100^2 \times 16) = 6 \times 10\,000 \times 16 = 960\,000$, на feed-forward энкодер: $6 \times (100 \times 16 \times 512) = 6 \times 819\,200 = 4\,915\,200$, аналогично для декодера ещё $960\,000 + 4\,915\,200 = 5\,875\,200$, на cross-attention: $6 \times (100 \times 100 \times 16) = 6 \times 160\,000 = 960\,000$. Итого: $960\,000 + 4\,915\,200 + 960\,000 + 4\,915\,200 + 960\,000 = 12\,710\,400$
FLOPs.

Также оценим и сложность модели из работы \cite{ge2022gnn}. Автор использовал модель \textbf{STGCN} (Spatial-Temporal Graph Convolutional Network) со следующими параметрами:

\begin{itemize}
	\item Вход: временной ряд OD-матриц (размер [batch, channels, time, nodes])
	\item 2 ST-Conv блока: каждый содержит 2 временных сверточных слоя (temporal 1D conv) и 1 пространственный слой (graph convolution)
	\item Пространственная свертка: через полиномы Чебышева (порядка K, в работе использовали K=3)
	\item Размерность графа: $n$ — число узлов (в работе были 15, 24, 50 узлов), OD-матрица — $n \times n$
	\item Временное окно: $T$ (используем 100)
	\item Размер батча: $B$ (В работе использовали 100)
	
\end{itemize}

Проведем оценку вычислительной сложности разных этапов:

\begin{itemize}
\item Временной сверточный 1D слой (GLU)

Сложность: $O(B \times n^2 \times T \times C_{in} \times C_{out})$ (на практике — 2 слоя, размерность ядра $k_t=3$)

Тогда оценка в FLOPs составит : $2 \times B \times n^2 \times T \times C_{in} \times C_{out}$

\item Пространственная графовая свертка (через полиномы Чебышева)

Сложность: $O(B \times n^2 \times K \times C_{in} \times C_{out})$

FLOPs: $2 \times B \times n^2 \times K \times C_{in} \times C_{out}$
\end{itemize}
	
Тогда всего на инференс для одного батча: 2 ST-Conv блока $\times$ (2 temporal + 1 spatial):

\begin{equation}
	\begin{split}
	O(B \times n^2 \times T \times C_{in} \times C_{out} \times 4 + B \times n^2 \times K \times C_{in} \times C_{out} \times 2) = \\
	O(B \times n^2 \times (T+K) \times C_{out})
	\end{split}
\end{equation}

То есть на один шаг для n=15,  будет затрачено 2 блока $\times$ 2 слоя $\times 15^2 \times 100 \times 1 \times 64 = 4 \times 225 \times 10 \times 64 = 5\,760\,000$ FLOPs и 2 блока $\times$ 1 слой $\times 15^2 \times 3 \times 1 \times 64 = 2 \times 225 \times 3 \times 64 = 86\,400$ FLOPs. Всего выходит: $\sim5\,846\,400$ FLOPs


Оценим также объем весов и параметров. Временная свертка - 2 блока $\times$ 2 слоя $\times (C_{in} \times C_{out} \times k_{t}) = 4 \times 1 \times 64 \times 3 = 768$ параметров.
Для пространственной свертки -  2 блока $\times (K \times C_{in} \times C_{out}) = 2 \times 3 \times 1 \times 64 = 384$ параметра
Всего: 768 + 384 = 1152 весов, но на практике чуть больше из-за bias и output слоя, поэтому округлим до 2тыс.). Таким образом стоит подчеркнуть, что количество параметров: зависит от числа каналов, фильтров, порядка полиномов Чебышева и числа слоев.

%2.4. Пример (n=15, T=10, K=3, C_in=1, C_out=64, B=1)
%Temporal: 2 блока \times 2 слоя \times 15² \times 10 \times 1 \times 64 = 4 \times 225 \times 10 \times 64 = 576,000 FLOPs
%Spatial: 2 блока \times 1 слой \times 15² \times 3 \times 1 \times 64 = 2 \times 225 \times 3 \times 64 = 86,400 FLOPs
%Всего: ~662,400 FLOPs (на один OD-ряд, n=15)
%Для n=50: (50²=2500) → ~7,360,000 FLOPs


Итоги проведенного анализа приведены в Таблицах \ref{tab:compare} и \ref{tab:compare2}.


\begin{table}[ht]
	\centering
	\caption{Сравнение вычислительной сложности, FLOPs и памяти моделей}
	\begin{tabular}{|c|c|>{\centering\arraybackslash}p{4.5cm}|c|}
		\hline
		\textbf{Модель} & \textbf{Параметров} & \textbf{Сложность} & \textbf{FLOPs} \\
		\hline
		LSTM & 33\,000 & $O((H+L) \times S^2 \times N)$ & $\sim819\,200$ \\ 
		\hline
		Transformer & 78\,000 & $O(N(H^2+L^2)S + N(H+L)S d_{ff})$ & $\sim12\,710\,400$ \\
		\hline
		STGCN (GNN) & $\sim2\,000$ & $O(B \times n^2 \times (T+K) \times C_{out})$ & $\sim5\,846\,400$ (n=15) \\ 
		\hline
		\textbf{ARIMA} & 5 & $O(T \times F)$ & $\sim800$ \\
		\hline
	\end{tabular}
	\label{tab:compare}
\end{table}




\begin{table}[ht]
	\centering
	\caption{Сравнение метрик рассматриваемых подходов}
	\begin{tabular}{|c|c|c|c|c|}
		\hline
		Метрика & GNN\cite{ge2022gnn} & Transformer\cite{lstm_trans} & LSTM\cite{lstm_trans} & ARIMA \\ 
		\hline
		MAE & 0,0023 & 0,1646 & 0,1833 & 0,1178  \\
		\hline
		RMSE & 0,0040 & - & - & 0,1493  \\ 
		\hline
		MSE & - & - & - & 0,0362  \\ 
		\hline
		MAPE & - & - & - & 2,4573  \\ 
		\hline
		\multicolumn{5}{|p{14cm}|}{\textit{Примечание}: прочерки означают, что авторы не предоставляют значения указанных метрик в своих работах.} \\
		\hline
	\end{tabular}
	\label{tab:compare2}
\end{table}

По итогам анализа можно подчеркнуть, что разработанный подход может демонстрировать чуть менее точные оценки, но при этом заметно менее вычислительно затратен, нежели методы, решающие аналогичные задачи.


Также стоит отметить, что для оценки точности использовались данные только <<неаномального>> кластера. В силу случайных изменений в аномальных кластерах ошибка оценки времени отклика может достигать RMSE = 0,5. Таким образом, исходя из баланса кластеров можно оценить ошибку на всех устройства как:

\begin{equation}
	\label{rmse_extim}
	RMSE_{all} \approx \sqrt(0,68 \times 0,15^2 + 0,32 \times 0,5^2 ) \approx 0,3
\end{equation}
Таким образом, фильтрация аномальных устройств обеспечивает снижение RMSE до 50\%.


\section{Выводы по результатам 4 главы}\label{subsec:ch4/sect5}

%Анализ времени отклика в различных регионах страны показал, что региональные особенности работы интернет-сети оказывают существенное влияние на скорость и надежность передачи данных. Ожидаемо устройства в близлежащих регионах показывают высокое быстродействие, однако детальный анализ показал наличие аномальных, выбивающихся из общей тенденции устройств, чье время отклика существенно выше ожидаемого. Анализ полученных временных рядов явно показал наличие сезонной компоненты в данных с большей части устройств с периодом, примерно равным 24 часам. Разделение временных рядов по паттернам поведения показало, что перед решением задачи предсказания быстродействия необходимо провести сегментирование наблюдаемого временного ряда, а также соотнести его с кластером схожих рядов с целью более точного предсказания.  Анализ моделей, решающих схожие задачи анализа времени отклика, показал что предлагаемый подход на 3-4 порядка быстрее (в терминах FLOPs) проанализированных работ.

%Применение результатов исследования может быть полезно для интернет-провайдеров и операторов связи, позволяя им оптимизировать свои сети и улучшать качество услуг для конечных пользователей.

Анализ динамических характеристик времени отклика в различных географических регионах Российской Федерации выявил значимое влияние региональных особенностей инфраструктуры телекоммуникационных сетей на показатели скорости и устойчивости передачи данных. Как и ожидалось, устройства, расположенные в непосредственной географической близости к точкам измерения, демонстрируют более высокие значения быстродействия. Тем не менее, детальное исследование выявило наличие аномальных наблюдений — устройств, чьё время отклика существенно превышает ожидаемые значения (выходит за 1,5 интерквартильных размаха от 1 и 3 квартиля остатков модели). Это отличается от общей тенденции и их объем составляет менее 10\% от наблюдений. Это указывает на возможные локальные сетевые узкие места (аномалии типа <<бутылочное горлышко>>) или особенности маршрутизации трафика.

При анализе временных рядов, соответствующих показателям rtt, была обнаружена выраженная сезонная компонента с периодом, близким к 24 часам, что свидетельствует о суточной цикличности нагрузки на сеть и, вероятно, связано с паттернами пользовательской активности. Кластеризация временных рядов по характеру динамики позволила выделить группы схожих поведенческих паттернов. Полученные результаты указывают на необходимость выполнения предварительного сегментирования временных рядов и их кластеризации перед построением моделей прогнозирования, что повышает точность предсказаний за счёт учёта гетерогенности сетевого трафика. Также рассмотренные методы кластеризации позволяют выделить аномальное поведение узлов во времени, которое отличается наличием иных шаблонов помимо суточной активности. 

Сравнительный анализ предложенного подхода с существующими методами решения задач прогнозирования времени отклика показал его существенное преимущество в вычислительной эффективности: предлагаемая модель требует на 3–4 порядка меньше операций с плавающей запятой (FLOPs), что делает её пригодной для масштабирования и применения в условиях ограниченных вычислительных ресурсов.

Практическая значимость полученных результатов заключается в возможности их использования телекоммуникационными операторами и интернет-провайдерами для диагностики сетевых аномалий, оптимизации топологии сетей передачи данных и повышения качества обслуживания конечных пользователей за счёт более точного прогнозирования и управления сетевыми задержками.


\clearpage
